% Funis user manual. Build with: .venv/bin/python docs/manual/make_manual.py
% The files under build/ are written by make_manual.py from the program itself.
\documentclass[fontsize=10.5pt,paper=a4,parskip=half,headings=normal,numbers=noenddot,
  open=any,bibliography=totoc,captions=tableheading]{scrreprt}

\usepackage[top=2.6cm,bottom=2.8cm,left=2.5cm,right=2.5cm,headsep=0.9cm,footskip=1.2cm]{geometry}
\usepackage{fontspec}
\usepackage{unicode-math}
\setmainfont{STIX Two Text}
\setmathfont{STIX Two Math}
\setsansfont{texgyreheros}[Extension=.otf,UprightFont=*-regular,BoldFont=*-bold,
  ItalicFont=*-italic,BoldItalicFont=*-bolditalic,Scale=0.92]
\setmonofont{Inconsolatazi4}[Extension=.otf,UprightFont=*-Regular,BoldFont=*-Bold,Scale=0.92]
\usepackage{newunicodechar}
\newunicodechar{≤}{\ensuremath{\leq}}
\newunicodechar{≥}{\ensuremath{\geq}}
\newunicodechar{×}{\ensuremath{\times}}
\newunicodechar{∞}{\ensuremath{\infty}}
\newunicodechar{≈}{\ensuremath{\approx}}
\newunicodechar{→}{\ensuremath{\rightarrow}}
\newfontfamily\symbolfont{STIX Two Math}
\newunicodechar{⌘}{{\symbolfont ⌘}}
\newunicodechar{⇧}{{\symbolfont ⇧}}
\newunicodechar{₀}{\textsubscript{0}}
\newunicodechar{₁}{\textsubscript{1}}
\newunicodechar{₂}{\textsubscript{2}}
\newunicodechar{₃}{\textsubscript{3}}
\newunicodechar{ₜ}{\textsubscript{t}}
\newunicodechar{ₙ}{\textsubscript{n}}
\newunicodechar{⁻}{\textsuperscript{−}}
\newunicodechar{⁺}{\textsuperscript{+}}

\usepackage[dvipsnames]{xcolor}
\definecolor{ink}{HTML}{243B53}
\definecolor{accent}{HTML}{1F6FEB}
\definecolor{muted}{HTML}{6B7A8C}
\definecolor{border}{HTML}{D9E0E8}
\definecolor{panel}{HTML}{F4F6F9}
\definecolor{reaction}{HTML}{A04460}
\definecolor{vertical}{HTML}{007C83}
\definecolor{horizontal}{HTML}{B57726}

\usepackage{graphicx}
\usepackage{float}
\usepackage{needspace}
\graphicspath{{build/figures/}{../../funis/assets/}}
\usepackage{tikz}
\usetikzlibrary{calc,positioning}
\usepackage{booktabs,tabularx,longtable,array,multirow,xltabular}
\usepackage{enumitem}
\setlist{itemsep=2pt,topsep=3pt,parsep=0pt}
\setlist[description]{font=\sffamily\bfseries\color{ink}}
\usepackage[font=small,labelfont={sf,bf},labelsep=period,skip=5pt]{caption}
\usepackage[most]{tcolorbox}
\usepackage{microtype}
\usepackage[hyphens]{url}
\usepackage[colorlinks=true,linkcolor=ink,urlcolor=accent,citecolor=ink,
  pdfauthor={Carlo Olivieri},pdfsubject={User manual of Funis},
  pdfkeywords={Funis, Linear Arch Static Analysis, LASA, form-finding, funicular, thrust line, masonry arches},
  bookmarksnumbered=true]{hyperref}
\usepackage{scrlayer-scrpage}

\input{build/generated/metadata.tex}
\hypersetup{pdftitle={Funis \FunisVersion{} User Manual}}

% ----------------------------------------------------------------- headings and page style
\addtokomafont{disposition}{\color{ink}}
\addtokomafont{chapter}{\Huge}
\renewcommand*{\chapterformat}{\mbox{\color{accent}\thechapter\enskip}}
\RedeclareSectionCommand[beforeskip=0pt,afterskip=1.4\baselineskip]{chapter}
\RedeclareSectionCommand[beforeskip=1.3\baselineskip,afterskip=0.5\baselineskip]{section}
\RedeclareSectionCommand[beforeskip=1.0\baselineskip,afterskip=0.3\baselineskip]{subsection}
\pagestyle{scrheadings}
\clearpairofpagestyles
\automark[chapter]{chapter}
\ihead{\sffamily\small\color{muted}\leftmark}
\ohead{\sffamily\small\color{muted}Funis \FunisVersion{} User Manual}
\cfoot*{\sffamily\small\color{muted}\pagemark}
\setkomafont{pageheadfoot}{\normalfont}
\renewcommand{\arraystretch}{1.18}
\setlength{\emergencystretch}{1.5em}
\renewcommand{\topfraction}{0.9}
\renewcommand{\bottomfraction}{0.8}
\renewcommand{\textfraction}{0.07}
\renewcommand{\floatpagefraction}{0.75}
\setcounter{topnumber}{3}
\setcounter{totalnumber}{5}
\numberwithin{figure}{chapter}
\numberwithin{table}{chapter}

% ----------------------------------------------------------------- notation and interface elements
\newcommand{\e}{\hat{\mathbf e}}
\newcommand{\x}{\mathbf x}
\newcommand{\F}{\mathbf F}
\newcommand{\G}{\mathbf G}
\newcommand{\R}{\mathbf R}
\newcommand{\C}{\mathbf C}
\newcommand{\Pv}{\mathbf P}
\newcommand{\T}{\hat{\mathbf T}}
\newcommand{\tv}{\hat{\mathbf t}}
\newcommand{\nv}{\hat{\mathbf n}}
\newcommand{\uv}{\mathbf u}
\newcommand{\pp}{\mathbf p}
\newcommand{\rr}{\mathbf r}
\newcommand{\dv}{\hat{\mathbf d}}
\newcommand{\mv}{\hat{\mathbf m}}
\newcommand{\menu}[1]{{\sffamily\bfseries\color{ink}#1}}
\newcommand{\ui}[1]{{\sffamily\color{ink}#1}}
\newcommand{\file}[1]{\texttt{#1}}
\newcommand{\kbd}[1]{\tikz[baseline=(k.base)]\node[draw=muted,rounded corners=1.5pt,inner xsep=2.5pt,
  inner ysep=1.2pt,font=\sffamily\footnotesize,fill=panel](k){#1};}
\newcommand{\badge}[1]{\tikz[baseline=(b.base)]\node[circle,fill=accent,text=white,inner sep=0.8pt,
  minimum size=12pt,font=\sffamily\bfseries\scriptsize](b){#1};}

% numbered callout on a screenshot: number, target x, y, badge x, y (fractions of the image)
\newcommand{\callout}[5]{\draw[accent,line width=0.6pt] (#4,#5) -- (#2,#3);
  \fill[accent] (#2,#3) circle (1.4pt);
  \node[circle,fill=accent,text=white,inner sep=0.8pt,minimum size=13pt,font=\sffamily\bfseries\scriptsize]
    at (#4,#5) {#1};}

\newtcolorbox{note}[1][Note]{enhanced,breakable,colback=panel,colframe=accent,boxrule=0pt,
  leftrule=2.5pt,arc=0pt,outer arc=0pt,left=8pt,right=8pt,top=5pt,bottom=5pt,
  fonttitle=\sffamily\bfseries\small,coltitle=accent,attach title to upper={\enskip},title={#1}}
\newtcolorbox{caution}[1][Caution]{enhanced,breakable,colback=reaction!6,colframe=reaction,boxrule=0pt,
  leftrule=2.5pt,arc=0pt,outer arc=0pt,left=8pt,right=8pt,top=5pt,bottom=5pt,
  fonttitle=\sffamily\bfseries\small,coltitle=reaction,attach title to upper={\enskip},title={#1}}

% screenshot with a thin frame: \shot[width]{file}
\newcommand{\shot}[2][\linewidth]{{\setlength{\fboxsep}{0pt}\setlength{\fboxrule}{0.4pt}%
  \color{border}\fbox{\includegraphics[width=\dimexpr#1-0.8pt\relax]{#2}}}}

% table of the fields of an input section
\newenvironment{fields}{\par\needspace{6\baselineskip}\small\setlength{\LTpre}{2pt}\setlength{\LTpost}{2pt}%
  \xltabular{\linewidth}{@{}>{\leavevmode\color{ink}\sffamily\raggedright\arraybackslash}p{0.33\linewidth}>{\raggedright\arraybackslash}p{0.11\linewidth}X@{}}
  \toprule \normalfont\bfseries Field & \bfseries Unit & \bfseries Meaning\\ \midrule\endhead
  \bottomrule\endfoot}
  {\endxltabular}

% summary of a worked example, generated from the program
\newcommand{\summarytable}[1]{{\par\needspace{10\baselineskip}\small\setlength{\LTpre}{2pt}\setlength{\LTpost}{2pt}%
  \begin{xltabular}{\linewidth}{@{}>{\leavevmode\color{ink}\sffamily\raggedright\arraybackslash}p{0.34\linewidth}>{\raggedright\arraybackslash}X@{}}
  \toprule \normalfont\bfseries Quantity & \bfseries Value\\ \midrule\endhead
  \bottomrule\endfoot #1 \end{xltabular}\par}}

\input{build/generated/summaries.tex}
\input{build/generated/examples.tex}

\begin{document}

% ----------------------------------------------------------------- title page
\begin{titlepage}
\begin{tikzpicture}[remember picture,overlay]
  \fill[ink] (current page.north west) rectangle ([yshift=-9.2cm]current page.north east);
  \fill[accent] ([yshift=-9.2cm]current page.north west) rectangle ([yshift=-9.45cm]current page.north east);
  \node[anchor=north west,inner sep=0] at ([xshift=2.5cm,yshift=-2.4cm]current page.north west)
    {\includegraphics[width=2.6cm]{icon.png}};
  \node[anchor=north west,text=white,font=\sffamily\bfseries\fontsize{54}{58}\selectfont]
    at ([xshift=5.6cm,yshift=-2.25cm]current page.north west) {Funis};
  \node[anchor=north west,text=white,font=\sffamily\fontsize{22}{26}\selectfont]
    at ([xshift=5.65cm,yshift=-4.25cm]current page.north west) {User Manual};
  \node[anchor=north west,text=white!80!ink,font=\sffamily\large,text width=13.5cm,align=left]
    at ([xshift=2.5cm,yshift=-6.4cm]current page.north west)
    {Equilibrium and form-finding of spatial compression paths:\\
     the analytical Linear Arch Static Analysis (LASA), plane funiculars\\
     and the thrust line assessment of arches};
  \node[anchor=north,inner sep=0] at ([yshift=-10.6cm]current page.north)
    {\shot[15.6cm]{window_overview.png}};
  \node[anchor=south west,font=\sffamily,text=ink,text width=16cm,align=left]
    at ([xshift=2.5cm,yshift=2.3cm]current page.south west)
    {{\Large\bfseries Version \FunisVersion}\\[3pt]
     {\color{muted}\ManualDate}\\[12pt]
     {\large \FunisAuthor}\\[2pt]
     {\small\color{muted}\FunisAffiliation}\\[8pt]
     {\small\color{muted}DOI \href{https://doi.org/\FunisDOI}{\FunisDOI}\quad\textbullet\quad
      \href{\FunisRepository}{\nolinkurl{\FunisRepository}}}};
\end{tikzpicture}
\end{titlepage}

% ----------------------------------------------------------------- imprint
\thispagestyle{empty}
\vspace*{\fill}
{\small\sloppy
\noindent\textsf{\bfseries\color{ink} Funis \FunisVersion{} User Manual}\\
\ManualDate

\medskip
\noindent This manual describes version \FunisVersion{} of Funis, the desktop application for Windows and macOS
of the analytical Linear Arch Static Analysis of spatial compression paths, with the plane funicular and
the thrust line assessment of masonry arches. It is rebuilt for every release of the program, and the
edition matching each version is kept in the repository as \file{docs/Funis\_User\_Manual.pdf}.

\medskip
\noindent\textsf{\bfseries\color{ink} Author.} \FunisAuthor, \FunisAffiliation.
ORCID \href{\FunisORCID}{\nolinkurl{\FunisORCID}}.

\noindent\textsf{\bfseries\color{ink} Repository.} \url{\FunisRepository}

\noindent\textsf{\bfseries\color{ink} Archive.} Every release is archived on Zenodo. The concept DOI
\href{https://doi.org/\FunisDOI}{\FunisDOI} refers to all versions and resolves to the latest one.

\noindent\textsf{\bfseries\color{ink} How to cite.} \FunisAuthor, \emph{Funis: equilibrium and
form-finding of spatial compression paths}, version \FunisVersion, software,
\href{https://doi.org/\FunisDOI}{doi:\FunisDOI}. The file \file{CITATION.cff} of the repository gives the
same data in machine-readable form.

\medskip
\noindent\textsf{\bfseries\color{ink} Licence.} Copyright \textcopyright{} 2026 \FunisAuthor. Funis is
released under the MIT License, given in the file \file{LICENSE} of the repository. This manual is part
of the repository and is distributed under the same licence. The installers also contain third-party
libraries under their own licences, among them Qt for Python (PySide6, LGPL-3.0), NumPy, SciPy and
Matplotlib.

\medskip
\noindent\textsf{\bfseries\color{ink} Disclaimer.} The software is provided ``as is'', without warranty
of any kind, as stated in the licence. Funis computes equilibrium states under the hypotheses described
in this manual. Whether those hypotheses represent a given structure, and every decision taken on the
basis of the results, remain the responsibility of the user.
}
\clearpage

\pdfbookmark[0]{Contents}{toc}
\setcounter{tocdepth}{1}
\tableofcontents

% =================================================================
\chapter{Introduction}\label{ch:intro}

\section{What Funis does}

Funis computes the equilibrium of continuous compression paths, that is, of lines that carry their loads
by an axial force alone. In space it follows the analytical Linear Arch Static Analysis (LASA) of
C.~Olivieri~\cite{article}: a path is either prescribed, and the program finds the axial force and the
lateral reactions that keep it in equilibrium, or unknown, and the program finds the curve from the
loads and the force at its start. In the vertical plane it draws the funicular line of given loads, the
catenary and the parabola among them, and assesses masonry arches of finite thickness by their thrust
lines.

The program is a desktop application for Windows and macOS. It needs no other software: the installers
contain everything, Python included. A model is entered in a panel of input sections, solved in a
fraction of a second, and read in a summary, in drawings to scale, in graphs along the line, in contact
diagrams, in a rotatable spatial view and in tables. Every drawing and table can be exported, and a
self-contained HTML report collects the whole analysis.

\section{The six tasks}

Table~\ref{tab:tasks} lists the tasks. The first four act on spatial paths and follow the article. The
last two act in the vertical plane.

\begin{table}[ht]
\caption{Tasks of Funis.}\label{tab:tasks}
\small
\begin{tabularx}{\linewidth}{@{}>{\sffamily\color{ink}\raggedright\arraybackslash}p{0.27\linewidth}X>{\raggedleft\arraybackslash}p{0.08\linewidth}@{}}
\toprule
\normalfont\bfseries Task & \bfseries Question answered & \bfseries Chapter\\ \midrule
Supported path (assessment) & A path of given shape leans on a lateral support, such as the wall of a
stair well. Which axial force and which reactions keep it in equilibrium, and can the support supply
them within a friction cone? & \ref{ch:supported}\\
Path without reactions & Is a path of given shape an axial equilibrium of given loads with no reaction
at all, and which loads would make it one? & \ref{ch:free}\\
Unknown geometry (form-finding) & Which curve carries given loads in pure compression from a given start
point and initial force, or between two given supports? & \ref{ch:form}\\
Node under a concentrated force & How does a concentrated force change the force and the direction of
the line at a node? & \ref{ch:node}\\
Plane funicular (catenary) & Which plane line carries given vertical and horizontal loads in compression
between two supports? & \ref{ch:funicular}\\
Arch assessment (thrust line) & Does a thrust line fit within a masonry arch, which thrusts are
admissible, and how far is the arch from collapse? & \ref{ch:arch}\\
\bottomrule
\end{tabularx}
\end{table}

\section{Scope and relation to the article}

The equations of the spatial tasks are those of the article~\cite{article}, and the cases of the article
are available in the \menu{Examples} menu, where figure and section numbers refer to it. The plane
funicular and the thrust line assessment of arches are tasks of the application. They rest on classical
statics of plane funiculars and on Heyman's safe theorem for masonry arches~\cite{heyman1966,heyman1969},
and they are not results of the article. Chapter~\ref{ch:limits} lists the hypotheses of every task.

\section{About this manual}

Chapters~\ref{ch:install} to~\ref{ch:conventions} cover installation, the interface and the conventions
common to every task. Chapters~\ref{ch:supported} to~\ref{ch:arch} describe the tasks one by one: the
input sections with their fields, the method, the results and a worked example taken from the
\menu{Examples} menu. Chapters~\ref{ch:checks} to~\ref{ch:limits} cover the checks of every result, the
exports, the catalogue of examples, the messages of the program and the limitations.
Appendix~\ref{app:quick} collects the keyboard shortcuts, the command line and the notation.

Interface elements are set in a sans-serif face: menu commands as \menu{File › Export › Report
(HTML)…}, fields and buttons as \ui{Preload T = −S(0)}, keys as \kbd{F5}. File names are set as
\file{project.funis}. The screenshots were taken on macOS. On Windows the window has the same layout,
with the fonts of the system.

The values printed in the worked examples and the screenshots are produced by the program itself when
this manual is built, so they are those that version~\FunisVersion{} gives.

% =================================================================
\chapter{Installation}\label{ch:install}

\section{Requirements}

The installers of every release are published on the release page of the repository,
\url{\FunisRepository/releases/latest}:
\begin{itemize}
\item \file{Funis-<version>-Windows-Setup.exe} for 64-bit Windows 10 or 11,
\item \file{Funis-<version>-macOS-arm64.pkg} and \file{Funis-<version>-macOS-arm64.dmg} for Macs with
Apple silicon.
\end{itemize}
On any other system, including Macs with an Intel processor and Linux, Funis runs from its source code
with Python~3.9 or later (Section~\ref{sec:source}). No licence key, account or network connection is
needed.

\section{Windows}

\begin{enumerate}
\item Download \file{Funis-<version>-Windows-Setup.exe} and run it.
\item The installer is not signed, so Windows SmartScreen may show the message \emph{Windows protected
your PC}. Choose \ui{More info} and then \ui{Run anyway}.
\item Follow the wizard. Funis is installed in \file{Program Files\textbackslash Funis}, with an entry
in the Start menu, an optional icon on the desktop and the association of the \file{.funis} files, so
that a double click on a project opens it in Funis.
\end{enumerate}
To uninstall Funis use \ui{Uninstall Funis} in the Start menu, or \ui{Settings › Apps}.

\section{macOS}

\begin{enumerate}
\item Download \file{Funis-<version>-macOS-arm64.pkg} and open it. The installer copies
\file{Funis.app} into \file{Applications}. Alternatively, open the \file{.dmg} file and drag
\file{Funis.app} into \file{Applications}.
\item The application is not notarised. At the first launch, macOS may refuse to open it. Open it with a
right click on \file{Funis.app} and \ui{Open}, or allow it in \ui{System Settings › Privacy \&
Security} with \ui{Open Anyway}. Later launches need no confirmation.
\end{enumerate}
To uninstall Funis move \file{Funis.app} from \file{Applications} to the Bin. On macOS, open project
files with \menu{File › Open…}.

\section{Running from the source code}\label{sec:source}

With Python~3.9 or later, in the folder of the repository:
\begin{verbatim}
python3 -m venv .venv
.venv/bin/python -m pip install -r requirements.txt
.venv/bin/python run_funis.py          # or: .venv/bin/python -m funis
\end{verbatim}
On Windows use \file{.venv\textbackslash Scripts\textbackslash python} instead of
\file{.venv/bin/python}. The dependencies are NumPy, SciPy, Matplotlib, PySide6, openpyxl, ezdxf and
Pillow. The test suite of the repository runs with
\begin{verbatim}
.venv/bin/python -m pip install pytest
.venv/bin/python -m pytest -q
\end{verbatim}

\section{Checking an installation}\label{sec:selftest}

The option \file{-{}-selftest} opens the window without showing it, solves every example and exits with
code~0 when all of them solve:
\begin{verbatim}
"C:\Program Files\Funis\Funis.exe" --selftest C:\Temp\funis.png      (Windows)
/Applications/Funis.app/Contents/MacOS/Funis --selftest ~/funis.png    (macOS)
\end{verbatim}
The optional image receives a screenshot of the window, and a report with the same name and the
extension \file{.log} is written next to it, for instance \emph{selftest: \NumberOfExamples{} examples,
0 failures}. The build scripts of the installers run the same test on the packaged program.

\section{Settings}

Funis remembers the size and position of the window, the state of \menu{Run automatically} and the last
folder used in a file dialog. They are stored with the user's preferences by Qt and contain no project
data.

% =================================================================
\chapter{Getting started}\label{ch:start}

\section{The main window}

At start, Funis opens the circular benchmark of the article (Figs.~11 and~12) and solves it.
Fig.~\ref{fig:window} shows the window. The \ui{Model} panel on the left holds the input, the results
fill the rest of the window, and every change of the input is solved again after a short pause.

\begin{figure}[ht]
\centering
\makebox[\linewidth]{%
\begin{tikzpicture}
  \node[anchor=south west,inner sep=0] (img) at (0,0) {\shot[\linewidth]{window_overview.png}};
  \begin{scope}[x={(img.south east)},y={(img.north west)}]
    \input{build/generated/callouts.tex}
  \end{scope}
\end{tikzpicture}}
\caption{The main window with the circular benchmark of the article, tab \ui{Views}.}\label{fig:window}
\end{figure}

\begin{description}[leftmargin=0pt,labelsep=0.5em]
\item[\badge{1} Menu bar.] \menu{File} (projects, examples, exports), \menu{Analysis} (run) and
\menu{Help} (user guide, about).
\item[\badge{2} Tool bar.] The most frequent commands: \ui{New project}, \ui{Open…}, \ui{Save},
\ui{Examples}, \ui{Run}, \ui{Run automatically}, \ui{Export} and \ui{User guide}.
\item[\badge{3} Problem.] The name of the project and the \ui{Task}. The task decides which input
sections appear below.
\item[\badge{4} Input sections.] The data of the task, in sections that scroll vertically: geometry,
loads, point actions, support, contact, output and so on (Section~\ref{sec:panel}).
\item[\badge{5} Result tabs.] \ui{Summary}, \ui{Views}, \ui{Graphs}, \ui{Contact} (\ui{Admissible
thrusts} for an arch), \ui{Spatial view} and \ui{Tables}. A tab without content for the task is
disabled.
\item[\badge{6} Selector.] \ui{All diagrams} of the tab, or one of them enlarged.
\item[\badge{7} Navigation.] The Matplotlib tool bar of the drawing: reset, back, forward, pan, zoom,
margins, axis options and save.
\item[\badge{8} Drawings.] Views to scale, with the labels placed clear of the lines and of each other.
\item[\badge{9} Notes.] The scales of the arrows and the meaning of dashed lines in the current
drawings.
\item[\badge{10} Status bar.] The outcome of the last run, for instance \emph{Supported path
(assessment) solved, 1 warning}, the messages of the exports and the version of the program.
\end{description}

\section{A first analysis}\label{sec:first}

The following steps use the benchmark that opens at start. They take a few minutes and touch every
part of the interface.

\begin{enumerate}
\item \textbf{Read the summary.} Open the \ui{Summary} tab. The thrust $S$ grows from $-8$~kN at the
start to about $-33$~kN at the end, the line is in compression everywhere, and the contact margin is
positive: the wall can supply the required reactions with the friction coefficient $\mu=0.6$.
\item \textbf{Change a load.} In \ui{Distributed loads}, raise the coefficient $a_0$ of the horizontal
load $h$ from 2 to 4 kN/m. After a short pause the analysis runs again. The summary now shows a warning
in colour: the required reaction leaves the contact cone. The line is still in compression, but the
assumed wall cannot supply the reaction (Section~\ref{sec:contact}).
\item \textbf{Find the admissible preloads.} The summary line \ui{Admissible preloads at λ = 1} gives
the interval of the preload $T$ for which the contact cone is satisfied with the new load, if any. Open
the \ui{Contact} tab and select \ui{Admissible preload and load multiplier} in the selector: the shaded
region contains every admissible pair $(T,\lambda)$, and the marker shows the model.
\item \textbf{Look at the line in space.} Open \ui{Spatial view} and drag the drawing to rotate it.
\item \textbf{Export a report.} Choose \menu{File › Export › Report (HTML)…}, and open the file in a web
browser. It holds the summary, every drawing, every table and the input.
\item \textbf{Save the project.} Choose \menu{File › Save} (\kbd{Ctrl+S}, \kbd{⌘S} on macOS). The project
is written as a \file{.funis} file that \menu{File › Open…} reads back.
\end{enumerate}

\section{Projects}

A project holds the whole input of a task: its name, the task and the data of every section. It does not
hold results, which are computed again when the project is opened.

\begin{description}[leftmargin=0pt,labelsep=0.5em]
\item[New project] (\kbd{Ctrl+N}) clears the input to the defaults of the program.
\item[Open…] (\kbd{Ctrl+O}) reads a \file{.funis} file, or a \file{.json} file with the same content.
\item[Save] (\kbd{Ctrl+S}) and \textsf{\bfseries\color{ink}Save as…} (\kbd{Ctrl+Shift+S}) write the
project. The title bar shows the file name, with an asterisk when there are unsaved changes.
\item[Examples] load one of the \NumberOfExamples{} examples of Chapter~\ref{ch:examples}. An example
is an ordinary project: modify it and save it under a new name.
\end{description}
Before a command replaces a project with unsaved changes, and before the window closes, Funis asks
whether to discard them. On macOS the shortcuts use \kbd{⌘} in place of \kbd{Ctrl}.

% =================================================================
\chapter{Working with the interface}\label{ch:interface}

\section{Menus and tool bar}

\begin{table}[ht]
\caption{Commands of the menus. The tool bar repeats the most frequent ones.}\label{tab:menus}
\small
\begin{tabularx}{\linewidth}{@{}>{\sffamily\color{ink}\raggedright\arraybackslash}p{0.34\linewidth}>{\raggedright\arraybackslash}p{0.17\linewidth}X@{}}
\toprule
\normalfont\bfseries Command & \bfseries Shortcut & \bfseries Action\\ \midrule
File › New project & Ctrl+N & Clear the input to the defaults.\\
File › Open… & Ctrl+O & Open a project (\file{.funis} or \file{.json}).\\
File › Examples & & Load an example (Chapter~\ref{ch:examples}).\\
File › Save & Ctrl+S & Save the project.\\
File › Save as… & Ctrl+Shift+S & Save the project under a new name.\\
File › Export › Drawings (PDF, SVG, PNG)… & & Write the selected drawings (Section~\ref{sec:exdrawings}).\\
File › Export › Drawings as DXF… & & Write the views to scale and the spatial path to one DXF file.\\
File › Export › Tables as CSV… & & Write every table to one CSV file.\\
File › Export › Tables as Excel… & & Write every table to a sheet of one workbook.\\
File › Export › Results as JSON… & & Write the input, the summary, the tables and the messages.\\
File › Export › Report (HTML)… & & Write a self-contained report.\\
File › Quit & Alt+F4 & Close Funis (Alt+F4 is the window command of Windows).\\
Analysis › Run & F5 & Solve the model now.\\
Analysis › Run automatically & & Solve the model after every change (on by default).\\
Help › User guide & F1 & The short guide of the program, with the equations of every task.\\
Help › About Funis & & Version, licence and reference.\\
\bottomrule
\end{tabularx}
\end{table}

The shortcuts are those of Windows. On macOS \kbd{⌘} replaces \kbd{Ctrl}, \kbd{⌘?} opens the user
guide, and \menu{Quit} (\kbd{⌘Q}) is in the application menu (Section~\ref{app:keys}). The \ui{Examples} and \ui{Export} buttons of the tool bar open the same lists as the
menus.

\section{The Model panel}\label{sec:panel}

The panel opens with the \ui{Problem} box (Fig.~\ref{fig:problem}): a free \ui{Name}, printed as the
title of the summary and of the report, and the \ui{Task}. Below it, the sections of the task appear in
a fixed order (Table~\ref{tab:sections}), and the panel scrolls when they do not fit. The panel can be
detached from the window with the button at its top right and docked again by a double click on its
title.

\begin{figure}[H]
\centering\shot[0.6\linewidth]{section_problem.png}
\caption{The Problem box.}\label{fig:problem}
\end{figure}

\begin{table}[ht]
\caption{Input sections of every task, in the order of the panel.}\label{tab:sections}
\small
\begin{tabularx}{\linewidth}{@{}>{\sffamily\color{ink}\raggedright\arraybackslash}p{0.30\linewidth}X@{}}
\toprule
\normalfont\bfseries Task & \bfseries Sections\\ \midrule
Supported path (assessment) & Geometry of the path, Distributed loads, Point actions, Support of the
line, Contact with the wall, Output\\
Path without reactions & Geometry of the path, Path without reactions, Distributed loads and Point
actions (only to check given loads), Output\\
Unknown geometry (form-finding) & Free curve, Distributed loads, Point actions, Output\\
Node under a concentrated force & Node\\
Plane funicular (catenary) & Plane funicular, Loads on the line, Point loads, Output\\
Arch assessment (thrust line) & Arch, Fill and surcharge, Thrust line assessment, Point loads, Output\\
\bottomrule
\end{tabularx}
\end{table}

\section{Entering values}

\begin{description}[leftmargin=0pt,labelsep=0.5em]
\item[Numbers.] A numeric field shows its unit after the value. Type the value, or use the arrows and the
mouse wheel. Lengths are in metres, forces in kilonewtons, angles in degrees.
\item[Choices.] A drop-down list selects a variant, for instance the generator of the geometry or the
kind of horizontal load. The fields that the variant does not use are disabled or hidden.
\item[Polynomials.] A load or a thrust given as a polynomial has the four coefficients $a_0$ to $a_3$ of
$a_0+a_1s+a_2s^2+a_3s^3$, with $s$ in metres. A constant is entered in $a_0$ alone.
\item[Vectors.] A vector has one field per component, for instance $C_1$, $C_2$, $C_3$ of an initial
force.
\item[Tables.] Points, load tables and point actions are entered in tables. \ui{Add row} appends a copy
of the last row, \ui{Remove} deletes the selected rows (or the last one), and \ui{Fill column} copies
the value of the selected cell to its whole column. A cell accepts a decimal point or a decimal comma. A
cell that is not a number stops the run with a message naming its row and column. A table of loads is
interpolated linearly between its rows and held constant beyond its ends.
\end{description}

\section{Running the analysis}

With \menu{Analysis › Run automatically} on, every change is solved 0.35~s after the last keystroke, so
the results follow the input while it is typed. With the option off, or for a large number of samples,
press \kbd{F5} or \ui{Run}. The status bar reports the outcome:
\begin{itemize}
\item \emph{<task> solved}, possibly with the number of warnings, which the \ui{Summary} tab lists in
colour above the results,
\item \emph{Input error}, when a field cannot be read, or \emph{The task could not be solved}, with the
reason, when the data do not admit a solution of the task. The \ui{Summary} tab shows the message
(Chapter~\ref{ch:messages}).
\end{itemize}

\section{Reading the results}

\begin{description}[leftmargin=0pt,labelsep=0.5em]
\item[Summary.] The name, the task, the notes of an example, the warnings and a table of the main
results (Fig.~\ref{fig:summary}). The values can be selected and copied.
\item[Views.] Plans, elevations and other drawings to scale. Distributed loads and reactions are drawn
as arrows at the scale given in the notes below the drawing.
\item[Graphs.] Quantities along the line: forces, reactions, loads, residuals, curvature and torsion, the
eccentricity of an arch.
\item[Contact.] For a supported path with the contact check on: the contact margin, the contact cone,
the admissible region of preload and load multiplier, the admissible horizontal loads and the minimum
preload. For an arch the tab is named \ui{Admissible thrusts} and shows the range of the thrust against
the multiplier of the variable loads.
\item[Spatial view.] The line in space with its loads and reactions. Drag to rotate, drag with the right
button to zoom.
\item[Tables.] The values at evenly spaced stations, the balance checks, the summary and the tables of
the task (point actions, endpoint correction, joints, eccentricities, hinges, loads).
\end{description}

\begin{figure}[ht]
\centering\shot[0.86\linewidth]{results_summary.png}
\caption{The \ui{Summary} tab of the benchmark.}\label{fig:summary}
\end{figure}

In \ui{Views}, \ui{Graphs} and \ui{Contact}, the selector \ui{Show} displays \ui{All diagrams} of the tab
in a grid, or one of them enlarged. The tool bar at the right of the selector is that of Matplotlib:
\ui{Home} restores the full view, the arrows step through previous views, the cross pans, the magnifier
zooms to a rectangle, and the disk saves the visible figure as an image. Labels are placed again at
every redraw, clear of the lines, arrows and other labels, so they stay readable at any size of the
window.

\begin{figure}[p]
\centering
\shot[0.86\linewidth]{results_graphs.png}\\[8pt]
\shot[0.86\linewidth]{results_tables.png}
\caption{The \ui{Graphs} tab (top) and the \ui{Tables} tab (bottom) of the benchmark.}\label{fig:graphs}
\end{figure}

% =================================================================
\chapter{Conventions of the model}\label{ch:conventions}

\section{Units and signs}

\begin{itemize}
\item Lengths in metres, forces in kilonewtons, angles in degrees in the input and in radians in the
tables where the column says so. Unit weights in kN/m³, surcharges in kN/m².
\item Load densities in kN per metre of \emph{projected} length $s$, unless the input states another
measure (Section~\ref{sec:measure}).
\item The axes $\e_1$, $\e_2$ are horizontal and $\e_3$ points upward. A downward load has $p_3<0$.
\item Tension is positive. A compressed line has $S\le0$ and $N\le0$. In the plane tasks the thrust
$H=-F_1$ is positive in compression.
\item Reactions are forces applied \emph{to the line} by its supports.
\end{itemize}

\section{Geometry of a spatial path}

A spatial path is described by its projected arc length $s$, the length of its plan:
\begin{equation}
\x^*(s)=\x(s)+f(s)\,\e_3,\qquad \tv=\x_{,s},\qquad \nv=\e_3\times\tv,\qquad \tv_{,s}=\rho\,\nv,
\end{equation}
where $\x(s)$ is the plan, $f(s)$ the elevation, $\tv$ the projected unit tangent, $\nv$ its left normal
and $\rho$ the projected curvature. The grade is $f_{,s}$, and $J=\sqrt{1+f_{,s}^2}$ converts projected
length into actual length. The path runs in the direction of increasing $s$: its start is $s=0$ and its
end $s=L$.

\section{Forces}

The force on the positive cut of the line is
\begin{equation}
\F=S\,\uv=N\,\T,\qquad \uv=\tv+f_{,s}\,\e_3,\qquad S=\F\cdot\tv,\qquad N=S\,J,
\end{equation}
where $S$ is the thrust (the horizontal projection of the force along the plan tangent), $N$ the axial
force and $\T=\uv/J$ the unit tangent of the spatial path. The vertical component $V=F_3$ appears in the
tables of the supported path. Every local result satisfies the balance $\F_{,s}+\pp+\rr=\mathbf 0$, with
$\pp$ the load and $\rr$ the reaction per projected metre.

\section{Load measures}\label{sec:measure}

The vertical load is $p_3=-w$, where $w$ can be given:
\begin{itemize}
\item per projected metre, as a polynomial or a table in $s$, the measure of the article,
\item per metre of actual length, as a polynomial, which suits a self-weight. Then $p_3=-wJ$ per
projected metre.
\end{itemize}
The horizontal load can be absent, a profile $h(s)$ in a fixed direction $\beta$ measured from $\e_1$
toward $\e_2$ (a polynomial or a table), pseudo-static, $h=\gamma\,|p_3|$ along $\beta$, or given by two
polynomial components $p_1(s)$, $p_2(s)$. The multiplier $\lambda$ scales the horizontal load, so that
$\lambda=0$ removes it without changing its data. In the plane tasks the loads are per metre of the
horizontal abscissa $x$, except a weight per metre of actual length (Chapter~\ref{ch:funicular}).

\section{Point actions}

A point action is a concentrated force $\G=(G_1,G_2,G_3)$ at the abscissa $s$. The program adds a node
to the samples there, and every quantity has a value before and after the node. In the tables of
stations the node takes two rows. On a supported path a point action takes a horizontal point reaction,
and on a free curve the force jumps by $-\G$.

\section{Resolution}

The number of \ui{Samples along s} (2001 by default, between 101 and 20\,001) fixes the resolution of
the integrals along the line, computed by the Simpson rule. The \ui{Output} section
(Fig.~\ref{fig:output}) sets the abscissa of the graphs, the angle $\theta$ when the path has one or the
projected length $s$, and the number of \ui{Stations in the tables} (41 by default). The stations do not
affect the solution, only the rows of the tables and of the exports.

\begin{figure}[H]
\centering\shot[0.6\linewidth]{section_output.png}
\caption{The \ui{Output} section.}\label{fig:output}
\end{figure}

% =================================================================
\chapter{Supported path (assessment)}\label{ch:supported}

\section{Purpose}

The geometry of the path is prescribed, for instance the line of thrust of a helical stair or of a ramp,
and the path leans on a lateral support, such as the wall of the stair well, that supplies horizontal
reactions only ($r_3=0$). A preload $T=-S(0)$, the compressive end thrust at the start, completes the
data. The program finds the thrust, the axial force and the reactions that the support must supply, and
checks whether a wall with friction can supply them.

\section{Input}

\subsection*{Geometry of the path}

\begin{figure}[H]
\centering
\raisebox{-\height}{\shot[0.48\linewidth]{section_geometry.png}}\hfill
\raisebox{-\height}{\shot[0.48\linewidth]{section_ellipse_geometry.png}}
\caption{\ui{Geometry of the path}: circular helix (left) and elliptical helix (right).}\label{fig:geometry}
\end{figure}

\begin{fields}
Generator & & Circular helix, elliptical helix or smooth path through points.\\
Samples along s & & Resolution of the integrals along $s$.\\
Radius R & m & Circular helix $(R\cos\theta,R\sin\theta,c(\theta-\theta_0))$, with $s=R(\theta-\theta_0)$
and grade $c/R$.\\
Pitch c (rise per radian) & m/rad & Negative for a path that descends as $\theta$ grows.\\
Start angle θ₀, end angle θ₁ & ° & The angular extent of the path, with $\theta_1>\theta_0$.\\
Semi-axis a, b & m & Elliptical plan $(a\cos\theta,b\sin\theta)$.\\
Elevation & & $f=c(\theta-\theta_0)$, or $f=k_s s$ with a constant grade.\\
c [m/rad] or k\textsubscript{s} [–] & & The rise per radian or the constant grade.\\
Points x₁, x₂, x₃ & m & A cubic spline through the points, re-parametrized by the projected length. The plan
must not stop or turn back.\\
\end{fields}

\begin{figure}[H]
\centering\shot[0.6\linewidth]{section_points_geometry.png}
\caption{\ui{Geometry of the path}: smooth path through points.}\label{fig:points}
\end{figure}

The supported path needs a non-zero grade everywhere, because the thrust follows from the vertical
balance divided by the grade.

\subsection*{Distributed loads}

\begin{figure}[H]
\centering\shot[0.6\linewidth]{section_loads.png}
\caption{\ui{Distributed loads (per metre of projected length)}.}\label{fig:loads}
\end{figure}

\begin{fields}
Vertical load p₃ = −w & & $w(s)$ per projected metre, $w(s)$ per metre of actual length, or a table of
$w(s)$ per projected metre.\\
w = a₀ + a₁ s + a₂ s² + a₃ s³ & kN/m & Coefficients of $w$ (or the table \ui{s}, \ui{w}).\\
Horizontal load & & None, $h(s)$ in a fixed direction $\beta$, a table of $h(s)$, pseudo-static
$h=\gamma|p_3|$ along $\beta$, or the components $p_1(s)$, $p_2(s)$.\\
h = a₀ + … & kN/m & Coefficients of $h$.\\
Direction β & ° & Direction of the horizontal load, from $\e_1$ toward $\e_2$.\\
Coefficient γ & – & Ratio of the pseudo-static horizontal load to the vertical load.\\
p₁, p₂ = a₀ + … & kN/m & Components of the horizontal load.\\
Multiplier λ & – & Factor on the horizontal load.\\
\end{fields}

\subsection*{Point actions}

\begin{figure}[H]
\centering\shot[0.6\linewidth]{section_actions_points.png}
\caption{\ui{Point actions} with one concentrated force.}\label{fig:actions}
\end{figure}

Each row gives the abscissa \ui{s} and the components \ui{G₁}, \ui{G₂}, \ui{G₃} in kN of a concentrated
force.

\subsection*{Support of the line and contact with the wall}

\begin{figure}[H]
\centering
\raisebox{-\height}{\shot[0.48\linewidth]{section_support.png}}\hfill
\raisebox{-\height}{\shot[0.48\linewidth]{section_contact.png}}
\caption{\ui{Support of the line} (left) and \ui{Contact with the wall} (right).}\label{fig:support}
\end{figure}

\begin{fields}
Preload T = −S(0) & kN & Compressive end thrust supplied at the start of the path.\\
Wall & & On the right or on the left of the direction of travel.\\
Check the contact cone & & Switches on the contact check and the \ui{Contact} tab.\\
Friction coefficient μ & – & Friction between the line and the wall.\\
Width of the centred strip L\textsubscript{sw} & m & Width of the strip of contact.\\
Distance of the wall δ & m & Distance of the strip from the wall. 0 removes the strip limit.\\
Largest T, largest λ of the diagrams & kN, – & Window of the diagram of admissible pairs $(T,\lambda)$.
0 chooses it automatically.\\
\end{fields}

\section{Method}

Vertical balance gives the vertical component of the force, $V=C_3-P_3(s)$, with $\C=\F(0)$ and $P_3$
the primitive of the vertical loads. Hence
\begin{equation}
S=\frac{V}{f_{,s}},\qquad N=S\,J,\qquad
r_t=-S_{,s}-\pp_h\cdot\tv,\qquad r_n=-\rho\,S-\pp_h\cdot\nv,
\end{equation}
where $r_t$ and $r_n$ are the tangential and normal components of the horizontal reaction and $\pp_h$
the horizontal load. A horizontal load leaves the thrust unchanged and shifts the reaction by $-\pp_h$:
the dashed arrows of the plan show the reaction of the vertical loads alone.

A point action $\G$ at $s_j$, with no vertical support impulse, makes the thrust jump and requires a
horizontal point reaction:
\begin{equation}
[S]=-\frac{G_3}{f_{,s}},\qquad \R_h=\frac{G_3}{f_{,s}}\,\tv-\G_h.
\end{equation}

\section{Contact with the wall}\label{sec:contact}

The wall presses the line with the normal reaction $r_\nu=r_n$ for a wall on the right of the direction
of travel and $r_\nu=-r_n$ for a wall on the left. The reaction is admissible when it lies in the
contact cone
\begin{equation}
r_\nu\ge0,\qquad |r_t|\le\eta\,r_\nu,\qquad \eta=\min(\mu,\varphi),\qquad \varphi=\frac{L_{sw}}{2\delta},
\end{equation}
where $\varphi$ limits the tangential reaction that a centred strip of width $L_{sw}$ at the distance
$\delta$ of the wall can carry. With $\delta=0$ the strip limit is removed. The contact margin
$\min(r_\nu,\ \eta r_\nu-|r_t|)$ is non-negative exactly inside the cone.

Every reaction is affine in the preload $T$ and in the load multiplier $\lambda$. The program therefore
gives in closed form:
\begin{itemize}
\item the interval of admissible preloads at the given $\lambda$, and the side of the cone and the
station that govern it,
\item the interval of admissible multipliers at the given $T$, when there is a horizontal load,
\item the convex region of admissible pairs $(T,\lambda)$,
\item the admissible horizontal loads turned to every direction $\beta$ at the given preload,
\item the minimum preload against the size of the horizontal load, along its direction and along the
directions turned by 90° and 180°,
\item for a descending circular helix with constant loads per projected metre, a wall on the right and no
point actions, the exact minimum preload of Appendix~A of the article.
\end{itemize}

\begin{caution}
A violated contact cone means that the assumed wall cannot supply the required reaction. It does not
mean that the line loses compression, and the bounds on $T$ and $\lambda$ are not collapse loads.
\end{caution}

\section{Results}

\begin{description}[leftmargin=0pt,labelsep=0.5em]
\item[Summary.] Geometry, projected and actual length, drop, preload, ranges of $S$, $N$, $r_t$ and
$r_n$, end forces, compression, wall, global balance and, with the contact check, the cone, the least
margin, the admissible preloads and multipliers and the exact bound when it applies.
\item[Views.] \ui{Plan with loads and reactions} and \ui{Developed elevation}.
\item[Graphs.] \ui{Thrust and axial force} and \ui{Required reactions}.
\item[Contact.] \ui{Contact margin}, \ui{Contact cone}, \ui{Admissible preload and load multiplier},
\ui{Admissible horizontal loads at T = …} and \ui{Minimum preload along the load ray}.
\item[Tables.] \ui{Stations} ($s$, $\theta$, position, grade, $S$, $N$, $V$, $r_t$, $r_n$, loads and
margin), \ui{Point actions} (jump of $S$, point reaction and its cone margin), \ui{Balance} and
\ui{Summary}.
\end{description}

\begin{figure}[ht]
\centering
\includegraphics[width=0.485\linewidth]{sup_plan.pdf}\hfill\includegraphics[width=0.485\linewidth]{sup_cone.pdf}\\[4pt]
\includegraphics[width=0.485\linewidth]{sup_region.pdf}\hfill\includegraphics[width=0.485\linewidth]{sup_domain.pdf}
\caption{Drawings of the benchmark exported by Funis: plan with loads and reactions, contact cone,
admissible pairs of preload and load multiplier, admissible horizontal loads at the given preload.}
\label{fig:supdrawings}
\end{figure}

\section{Worked example: the circular benchmark}

Load \menu{Examples › \NameBenchmark}. The path is a complete descending turn of a circular helix of
radius 2~m and pitch $-1$~m/rad, under a vertical load of 1~kN per projected metre and a horizontal load
of 2~kN/m along $\e_1$, with a preload of 8~kN and a wall on the right. The contact cone has $\mu=0.6$
and a strip of 0.8~m at 0.5~m from the wall. The summary reads:

\summarytable{\SummaryBenchmark}

\medskip
The contact margin is positive, so the wall can supply the reactions. The admissible preloads start
at the exact bound of Appendix~A, and the horizontal load could grow by the factor printed as the upper
admissible multiplier before the cone is violated. The example \menu{\NameViolated} doubles the
horizontal load: the warning
\begin{itemize}\small\MessagesViolated\end{itemize}
appears above the summary, and Fig.~\ref{fig:violated} shows it.

\begin{figure}[ht]
\centering\shot[0.86\linewidth]{results_warning.png}
\caption{A warning in the \ui{Summary} tab: the reaction leaves the contact cone.}\label{fig:violated}
\end{figure}

% =================================================================
\chapter{Path without reactions}\label{ch:free}

\section{Purpose}

Without reactions the force along the line is $\F=\C-\Pv(s)$, with $\C$ the force at the start and
$\Pv$ the primitive of the distributed loads and of the point forces. A prescribed path is an axial
equilibrium only if $\F$ is parallel to $\uv$ everywhere. This task checks that condition for given
loads, or finds the loads that satisfy it.

\section{Input}

The sections \ui{Geometry of the path} and \ui{Output} are those of Chapter~\ref{ch:supported}. The
section \ui{Path without reactions} (Fig.~\ref{fig:rf}) selects the task, and \ui{Distributed loads}
and \ui{Point actions} appear only to check given loads.

\begin{figure}[H]
\centering\shot[0.6\linewidth]{section_reaction_free.png}
\caption{The section \ui{Path without reactions}.}\label{fig:rf}
\end{figure}

\begin{fields}
Task & & Check the loads below, loads in a fixed vertical plane, or loads for a prescribed thrust
$S(s)$.\\
Conserved component B = F · m & kN & For loads in a fixed vertical plane.\\
Plane direction β & ° & Direction $\dv(\beta)$ of the vertical plane of the loads.\\
S = a₀ + a₁ s + … & kN & The prescribed thrust.\\
\end{fields}

\section{Method}

\begin{description}[leftmargin=0pt,labelsep=0.5em]
\item[Check the loads.] The program finds the initial force $\C$ that makes $\F=\C-\Pv$ closest to the
tangent in the least-squares sense, and reports the largest transverse force $|\F-S\uv|$. A transverse
force that is not small, compared with the force itself, disproves the equilibrium for every $\C$: a
lateral support or a different geometry is needed.
\item[Loads in a fixed vertical plane.] Loads in the plane spanned by $\e_3$ and $\dv(\beta)$ conserve
the component $\F\cdot\mv=B$ along $\mv=\e_3\times\dv$, so that $S=B/(\tv\cdot\mv)$ and the compatible
loads follow from $\pp=-\F_{,s}$. Compression needs $\tv\cdot\mv$ of opposite sign to $B$, which
restricts the path to part of a half-turn.
\item[Loads for a prescribed thrust.] The loads are $\pp=-(S\uv)_{,s}$, in general with three
components.
\end{description}

\section{Results}

The summary gives the verdict or the compatible loads, the ranges of $S$ and $N$, the compression and the
global balance. The views are the plan with the horizontal loads and the developed elevation, and the
graphs show the forces, the loads and, for a check, \ui{Transverse force of the best fit}. A part of the
path in tension is reported with a warning.

\section{Worked example}

\menu{Examples › \NameNotAxial} perturbs the elevation of a funicular curve. The check finds:

\summarytable{\SummaryNotAxial}

\medskip
\noindent with the warning
\begin{itemize}\small\MessagesNotAxial\end{itemize}
The example that follows it in the menu solves the same check on the funicular itself, which passes.

\begin{figure}[ht]
\centering\includegraphics[width=0.485\linewidth]{rf_residual.pdf}\hfill\includegraphics[width=0.485\linewidth]{rf_loads.pdf}
\caption{Transverse force of the best fit for the perturbed path (left), and the compatible loads of
the exceptional helical arc of Fig.~8 of the article (right).}\label{fig:rfdrawings}
\end{figure}

% =================================================================
\chapter{Unknown geometry (form-finding)}\label{ch:form}

\section{Purpose}

The loads are given as functions of the projected length and the curve is unknown. Starting from a
point and from the initial force, the program integrates the curve that carries the loads in pure
compression. It can also find the initial force that takes the curve to a prescribed second support.

\section{Input}

\begin{figure}[H]
\centering\shot[0.6\linewidth]{section_form.png}
\caption{The section \ui{Free curve} with a prescribed second support.}\label{fig:form}
\end{figure}

\begin{fields}
Projected length L & m & Length of the plan of the curve.\\
Samples along s & & Resolution of the integration.\\
Start x*(0) & m & The first support.\\
Task & & Prescribed initial force, or prescribed second support.\\
Initial force C & kN & $\C=\F(0)$, for a prescribed initial force.\\
Second support x*(L) & m & The target, for a prescribed second support.\\
Estimate of C & kN & Starting value of the correction, for a prescribed second support.\\
\end{fields}
The sections \ui{Distributed loads} and \ui{Point actions} are those of Chapter~\ref{ch:supported}.

\section{Method}

The force is $\F=\C-\Pv(s)$ and the curve follows by quadrature from
\begin{equation}
\x^*_{,s}=-\frac{\F}{Q},\qquad Q=|\F_h|>0,\qquad S=-Q,\qquad N=-|\F|,
\end{equation}
where $\F_h$ is the horizontal part of the force. With a weight per metre of actual length the force and
the position are integrated together. With a prescribed second support, the initial force is corrected by
a damped Newton iteration with the exact endpoint Jacobian (finite differences for a weight per actual
length), starting from the estimate. Every iterate keeps $Q>0$.

The graphs give the curvature and the torsion of the curve,
\begin{equation}
\kappa=\frac{Q\,|\F\times\pp|}{|\F|^3},\qquad
\tau=-\frac{Q\,\F\cdot(\pp\times\pp_{,s})}{|\F\times\pp|^2}.
\end{equation}

\section{Results}

The summary gives the end point, the end force, the length of the curve, the ranges of $S$ and $N$, the
least horizontal force $Q$, the distance from the plane of $\C$ and $\pp(0)$ (zero for parallel loads)
and, for constant loads, the difference from the closed form of Appendix~B of the article. With a
prescribed second support it adds the target, the number of force updates with the final mismatch, and
the eigenvalues of the plan block $\mathbf K$ of the endpoint Jacobian. The views are the plan and the
developed elevation, with \ui{Elevations of the successive estimates} for a second support. The graphs
show the forces, the curvature and torsion, the distance from the plane and the endpoint correction.

\section{Worked example}

\menu{Examples › \NameSecondSupport} prescribes the second support of case C of the article and starts
from a deliberately wrong estimate:

\summarytable{\SummarySecondSupport}

\begin{figure}[ht]
\centering\includegraphics[width=0.485\linewidth]{ff_plan.pdf}\hfill\includegraphics[width=0.485\linewidth]{ff_elevation.pdf}\\[4pt]
\includegraphics[width=0.485\linewidth]{ff_iterates.pdf}\hfill\includegraphics[width=0.485\linewidth]{ff_newton.pdf}
\caption{Form-finding: plan and developed elevation of case C (top), and the successive estimates and
the endpoint mismatch of the correction for a prescribed second support (bottom).}\label{fig:ffdrawings}
\end{figure}

\begin{figure}[ht]
\centering\shot[0.86\linewidth]{results_spatial.png}
\caption{The \ui{Spatial view} of case C. Drag to rotate, drag with the right button to zoom.}
\label{fig:spatial}
\end{figure}

% =================================================================
\chapter{Node under a concentrated force}\label{ch:node}

\section{Purpose and input}

At a node where a concentrated force $\G$ and possibly a point reaction $\R$ act, the force of the line
jumps. The task gives the force after the node and the change of direction of the line.

\begin{figure}[H]
\centering\shot[0.6\linewidth]{section_node.png}
\caption{The section \ui{Node}.}\label{fig:node}
\end{figure}

\begin{fields}
Force before the node F⁻ & kN & The force of the line before the node.\\
Point force G & kN & The concentrated force.\\
Point reaction R & kN & An optional point reaction.\\
Length of the drawn members & m & Length of the members in the drawings.\\
\end{fields}

\section{Method and results}

\begin{equation}
\F^+=\F^--\G-\R,\qquad \T^\pm=-\frac{\F^\pm}{|\F^\pm|}.
\end{equation}
The summary gives $\F^+$, the axial forces $N^-$ and $N^+$, the tangents, the change of tangent in
degrees and the closure of the force polygon. The views are the node in plan, in elevation on
$x_1x_3$ and the force polygon, drawn in the plane of the forces. For \menu{Examples ›
\NameNode}:

\summarytable{\SummaryNode}

\begin{figure}[ht]
\centering\includegraphics[width=0.485\linewidth]{node_plan.pdf}\hfill\includegraphics[width=0.485\linewidth]{node_polygon.pdf}
\caption{The node of Fig.~7 of the article in plan, and its force polygon.}\label{fig:nodedrawings}
\end{figure}

% =================================================================
\chapter{Plane funicular}\label{ch:funicular}

\section{Purpose}

The line lies in the vertical plane of $\e_1$ and $\e_3$, between the supports $A=(0,0)$ and
$B=(L,z_B)$, and is described by its height $z=f(x)$ over the horizontal abscissa $x$. The program finds
the line that carries the loads in compression: the catenary under a weight per metre of actual length,
the parabola under a load per metre of $x$, the funicular polygon of point loads, and their
combinations with horizontal loads.

\section{Input}

\begin{figure}[H]
\centering
\raisebox{-\height}{\shot[0.48\linewidth]{section_funicular.png}}\hfill
\raisebox{-\height}{\shot[0.48\linewidth]{section_plane_loads.png}}
\caption{The sections \ui{Plane funicular} (left) and \ui{Loads on the line} (right).}\label{fig:pf}
\end{figure}

\begin{fields}
Span L & m & Horizontal distance of the supports.\\
Height z\textsubscript{B} of the second support & m & Height of $B$ above $A$.\\
Samples along x & & Resolution of the integration.\\
Line & & Through a point and the second support, with a prescribed thrust and the second support, or from
a prescribed initial force.\\
Point (x\textsubscript{c}, z\textsubscript{c}) & m & A point of the line, with $0<x_c<L$.\\
Thrust H\textsubscript{A} = −F₁(0) & kN & The thrust at $A$, positive in compression.\\
Initial force F(0) & kN & The force at $A$. The line then ends where the loads take it.\\
Vertical load p₃ = −w & kN/m & $w(x)$ per metre of actual length (catenary), per metre of $x$
(parabola), or a table per metre of $x$.\\
Horizontal load (along e₁) & kN/m & None, $p_1=h(x)$ per metre of $x$, or pseudo-static $p_1=\gamma w$
in the measure of $w$.\\
Point loads x, G₁, G₃ & m, kN & Concentrated forces on the line. $G_3<0$ acts downward.\\
\end{fields}

\section{Method}

Without reactions the force is $\F=\C-\Pv(x)$, with $\Pv$ the primitive of the loads and of the point
loads, and the line follows $f_{,x}=F_3/F_1$. Compression needs $F_1<0$ everywhere. The thrust is
$H=-F_1$, the axial force $N=-|\F|$, and the reactions are $\R_A=-\F(0)$ and $\R_B=\F(L)$. A load per
metre of actual length is multiplied by $J=|\F|/|F_1|$, and the force and the line are integrated
together. A horizontal load makes the thrust vary along the line, and where $F_1$ would reach zero the
line turns vertical and the program stops with a message. A point load makes the force jump by $-\G$ and
the line has a corner there.

Where a closed form exists the summary gives its largest distance from the integrated line: the parabola
(constant $w$ per metre of $x$), the catenary (constant $w$ per metre of actual length), the logarithmic
branch of Appendix~B of the article (constant $w$ and $h$ per metre of $x$) and the catenary with an
inclined axis (constant $w$ per actual length with the pseudo-static load).

\section{Results and worked example}

The summary gives the loads, the selected line, the initial force, the thrusts and reactions at $A$ and
$B$, the highest point, the length, the range of $N$, the height reached at $x=L$ against the target,
the closed form when it applies and the global balance. For \menu{Examples › \NameCatenary}:

\summarytable{\SummaryCatenary}

\begin{figure}[ht]
\centering\includegraphics[width=0.485\linewidth]{pf_funicular.pdf}\hfill\includegraphics[width=0.485\linewidth]{pf_polygon.pdf}
\caption{A catenary under its weight and a pseudo-static horizontal action (left) and a funicular
polygon of point loads (right), with the reactions.}\label{fig:pfdrawings}
\end{figure}

% =================================================================
\chapter{Arch assessment (thrust line)}\label{ch:arch}

\section{Purpose}

The task assesses a masonry arch of finite thickness by its thrust lines, under the classical hypotheses
of Heyman~\cite{heyman1966}: the masonry has no tensile strength, an unlimited compressive strength, and
sliding between voussoirs does not occur unless a friction coefficient is given. By the safe theorem, if
one thrust line lies within the arch under the given loads, the arch does not collapse under those
loads.

\section{Input}

\begin{figure}[H]
\centering
\raisebox{-\height}{\shot[0.48\linewidth]{section_arch.png}}\hfill
\raisebox{-\height}{\shot[0.48\linewidth]{section_fill.png}}
\caption{The sections \ui{Arch} (left) and \ui{Fill and surcharge} (right).}\label{fig:archsec}
\end{figure}

\begin{fields}
Shape of the centre line & & Circular (segmental or semicircular), pointed (two-centred), elliptical,
parabolic, catenary, or through points $(x,z)$.\\
Span L & m & Distance of the centres of the springings, at $(0,0)$ and $(L,0)$.\\
Rise f & m & Rise of the centre line.\\
Thickness at the crown t, at the springings & m & The thickness varies linearly along the centre line
between them. The faces are offset from the centre line along its normal.\\
Depth b & m & Depth normal to the plane, which multiplies every area.\\
Unit weight of the masonry & kN/m³ & \\
Fill above the extrados & & Switches on the fill.\\
Level x₃ of the fill, unit weight & m, kN/m³ & The fill loads the extrados vertically, column by column,
and has no strength of its own.\\
Horizontal action also on the fill & & Applies the horizontal action to the weight of the fill too.\\
Uniform surcharge q & kN/m² & Vertical load on the extrados.\\
\end{fields}

\begin{figure}[H]
\centering\shot[0.6\linewidth]{section_assessment.png}
\caption{The section \ui{Thrust line assessment}.}\label{fig:assessment}
\end{figure}

\begin{fields}
Joints & & Normal to the centre line (voussoir joints) or vertical sections.\\
Number of slices & & The arch is cut into slices by the joints. Each slice carries its weight at its
centroid.\\
Horizontal action λ & – & A fraction of the weights applied horizontally at the same centroids.\\
Direction of the horizontal action & & Along $+\e_1$ or $-\e_1$.\\
Collapse multiplier of & & The horizontal action, the point loads, or none.\\
Friction coefficient & – & With $\mu>0$ every joint also needs $|V|\le\mu|N|$. With 0 sliding is
excluded.\\
Line in the tables & & Minimum thrust, maximum thrust, or the line through three points.\\
e / t of the three points & – & Eccentricities at the springing $A$, at the crown and at the springing
$B$, from $-1/2$ (intrados) to $1/2$ (extrados).\\
Minimum thickness & & Computes the least thickness, with the weights recomputed.\\
\end{fields}
The section \ui{Point loads} gives concentrated forces $(x,G_1,G_3)$ on the extrados at the abscissa
$x$.

\section{Method}

A thrust line is the line of the centres of pressure of the resultant forces on the joints. The force on
joint $j$ is $\F_j=\C-\Pv_j$, with $\C$ the force at the springing $A$ and $\Pv_j$ the sum of the loads
of the slices before the joint, and the moment follows in the same way, so that every line is fixed by
three numbers. The line is admissible where its centre of pressure lies on every joint between the
intrados and the extrados and the force compresses the joint. These conditions are linear in the three
numbers, and the program solves them as linear programs:
\begin{description}[leftmargin=0pt,labelsep=0.5em]
\item[Minimum and maximum thrust.] The admissible lines with the least and the largest horizontal thrust
$H_A$, with the joints where they touch a face.
\item[Line through three points.] The line with the given eccentricities at $A$, at the crown and at
$B$. It is reported even when it leaves the arch, with a warning.
\item[Geometric safety factor.] The joints are shrunk about their centres until a line just fits, with
the same loads. The factor is the actual thickness over the thickness of this reduced
arch~\cite{heyman1969}. A factor below 1 means that no line lies within the arch.
\item[Collapse multiplier.] The largest multiplier of the horizontal action or of the point loads for
which an admissible line exists. The line touches the faces at the hinges of the mechanism. With four
hinges and no friction limit, the summary checks the multiplier by the virtual work of the mechanism
through them.
\item[Minimum thickness.] The least factor on every thickness for which a line fits, with the weights
recomputed for the thinner arch. For the semicircular arch under its own weight on normal joints this is
the classical $t/R=0.1075$, with hinges near 54.5° from the crown.
\item[Friction.] With a coefficient $\mu>0$ the conditions $|V|\le\mu|N|$ join the linear programs. With
0 the summary gives the friction that the selected line needs.
\end{description}

\section{Results}

\begin{description}[leftmargin=0pt,labelsep=0.5em]
\item[Views.] \ui{Thrust lines in the arch}, \ui{Loads and reactions} of the selected line,
\ui{Collapse} with the hinges of the mechanism, \ui{Geometric safety factor} with the reduced arch, and
\ui{Arch of minimum thickness}.
\item[Graphs.] \ui{Eccentricity of the centre of pressure} and \ui{Normal force at the joints}.
\item[Admissible thrusts.] The range of $H_A$ of the lines in the arch against the multiplier of the
variable loads. It closes at the collapse multiplier.
\item[Tables.] \ui{Joints} (position, angle, length, forces, $e/t$, moment and centre of pressure of
the selected line), \ui{Eccentricities}, \ui{Hinges}, \ui{Admissible thrusts}, \ui{Loads}, \ui{Balance}
and \ui{Summary}.
\end{description}

\begin{figure}[ht]
\centering\shot[0.86\linewidth]{results_arch.png}
\caption{The \ui{Views} tab of a semicircular arch under a horizontal action.}\label{fig:archviews}
\end{figure}

\section{Worked examples}

\menu{Examples › \NameArchHorizontal} computes the collapse multiplier of a horizontal action
proportional to the weight of a semicircular arch with $t/R=0.15$:

\summarytable{\SummaryArchHorizontal}

\begin{figure}[ht]
\centering\includegraphics[width=0.485\linewidth]{arch_collapse.pdf}\hfill\includegraphics[width=0.485\linewidth]{arch_domain.pdf}\\[4pt]
\includegraphics[width=0.485\linewidth]{arch_gsf.pdf}\hfill\includegraphics[width=0.485\linewidth]{arch_thickness.pdf}
\caption{Collapse mechanism under the horizontal action, admissible thrusts against the multiplier, the
reduced arch of the geometric safety factor, and the semicircular arch of minimum thickness.}\label{fig:archdrawings}
\end{figure}

\menu{Examples › \NameArchThickness} finds the least thickness of the semicircular arch under its own
weight:

\summarytable{\SummaryArchThickness}

% =================================================================
\chapter{Checks of every result}\label{ch:checks}

Every result carries independent checks in the \ui{Balance} table:
\begin{itemize}
\item the resultant force and moment of all actions on the line, end forces and point actions included,
which vanish for an equilibrium,
\item a finite-difference test of the local equations $\F_{,s}+\pp+\rr=\mathbf 0$.
\end{itemize}
The last line of the summary repeats the global balance. The values are normalised by the sum of the
magnitudes of the actions. A value near the rounding error of the computation, $10^{-12}$ or below,
confirms the computed state. A larger value either reflects the resolution, and then decreases when the
samples are increased, or belongs to data that the program reports as admitting no equilibrium, such
as a path that is not an axial equilibrium of its loads, whose best fit leaves a moment unbalanced. For
an arch the balance is that of the selected line, with the reactions at the centres of pressure of the
springings and the loads of all the slices.

The test suite of the repository checks the solvers against the closed forms and against the values
reported in the article, the minimum thickness of the semicircular arch against the classical value, the
examples, the files and the exports, and the clearance of the labels of every drawing.

% =================================================================
\chapter{Exporting results}\label{ch:export}

Every export needs a solved model. The commands are in \menu{File › Export} and in the \ui{Export}
button of the tool bar.

\section{Drawings}\label{sec:exdrawings}

\menu{Drawings (PDF, SVG, PNG)…} opens the dialog of Fig.~\ref{fig:export}. Select the drawings, the
format and the size:
\begin{description}[leftmargin=0pt,labelsep=0.5em]
\item[Fit to a width.] Every drawing fills the given width with its axes and grid.
\item[Views to scale, without axes.] Plans, elevations and diagrams are drawn to scale, as in a plate,
with a scale bar and the title: the \ui{Length scale} in cm/m applies to plans and elevations, the
\ui{Force scale} in cm/kN to diagrams in kN or kN/m. Graphs along the line are always fitted to the
width.
\item[Resolution.] The dots per inch of a PNG image.
\end{description}
Each drawing is written as \file{<base>\_<diagram>.<format>}, where \file{<diagram>} is the key of
Table~\ref{tab:keys}. The labels are placed for the size of the file, clear of the geometry.

\begin{figure}[ht]
\centering\shot[0.46\linewidth]{export_dialog.png}
\caption{The dialog \ui{Export drawings}.}\label{fig:export}
\end{figure}

\begin{table}[ht]
\caption{Keys of the drawings in the exported file names.}\label{tab:keys}
\small
\begin{tabularx}{\linewidth}{@{}>{\ttfamily}p{0.2\linewidth}X@{}}
\toprule
\normalfont\bfseries Key & \bfseries Drawing\\ \midrule
plan, elevation & Plan with loads and reactions (or of the curve), developed elevation\\
forces, reactions, loads & Thrust and axial force, required reactions, loads\\
residual & Transverse force of the best fit (path without reactions)\\
curvature, plane & Curvature and torsion, distance from the plane of $\C$ and $\pp(0)$\\
newton, iterates & Endpoint correction, elevations of the successive estimates\\
margin, cone, region & Contact margin, contact cone, admissible preload and load multiplier\\
domain, preload & Admissible horizontal loads, minimum preload along the load ray\\
node\_plan, node\_elevation, polygon & Node in plan and in elevation, force polygon\\
funicular & Plane funicular\\
arch\_lines, arch\_loads & Thrust lines in the arch, loads and reactions\\
arch\_collapse, arch\_gsf, arch\_thickness & Collapse, geometric safety factor, arch of minimum
thickness\\
eccentricity, arch\_forces, thrust\_domain & Eccentricity, normal force at the joints, admissible
thrusts\\
\bottomrule
\end{tabularx}
\end{table}

\section{DXF}

\menu{Drawings as DXF…} writes one DXF file (AutoCAD 2010 format) with every view to scale, side by side
from left to right, each with its title. Plans and elevations are in metres, one drawing unit per metre.
Diagrams in kN or kN/m are drawn at 0.1 drawing units per unit. Every view has its own layers, named
after the drawing and the role of the line (loads, reactions and so on), and the spatial path of a
spatial task is written as a 3D polyline on the layer \file{spatial\_path}. Graphs of functions are not
written.

\section{Tables}

\begin{description}[leftmargin=0pt,labelsep=0.5em]
\item[CSV.] One file with every table, each introduced by its title and followed by its note, in UTF-8.
\item[Excel.] One workbook with a sheet per table, the summary first, a styled header and frozen
headings.
\item[JSON.] The title, the input of the project, the summary, every table with its columns, rows and
note, the messages and, for a spatial path, the coordinates of its nodes.
\end{description}

\section{Report}

\menu{Report (HTML)…} writes one self-contained HTML file, with the images embedded: the name and task,
the notes and warnings, the summary, every drawing, every table and the complete input. It opens in any
web browser and can be printed to PDF from there.

\section{Project files}

A \file{.funis} file is a JSON text with the keys \file{application}, \file{version} (the version of
Funis that saved it), \file{saved} (date and time) and \file{project} (the input). Fields missing from
an older file take their default values and unknown fields are ignored, so a project saved by one
version opens in the next. Being a text file, a project can be compared, versioned and generated by a
script.

% =================================================================
\chapter{Examples}\label{ch:examples}

The \menu{Examples} menu holds \NumberOfExamples{} projects (Table~\ref{tab:examples}): the cases of the
article, with figure and section numbers referring to it, three reference cases of the application (a
catenary, a parabolic arch and a free curve with a point force), and the plane cases, funiculars checked
against their closed forms and arches of several shapes. The descriptions are the notes that each
example shows in the \ui{Summary} tab.

{\small
\setlength{\LTpre}{4pt}
\begin{longtable}{@{}>{\sffamily\color{ink}\raggedright\arraybackslash}p{0.31\linewidth}>{\raggedright\arraybackslash}p{0.65\linewidth}@{}}
\caption{The examples, grouped as in the menu.}\label{tab:examples}\\
\toprule
\normalfont\bfseries Example & \bfseries Description\\ \midrule
\endfirsthead
\toprule
\normalfont\bfseries Example & \bfseries Description\\ \midrule
\endhead
\bottomrule
\endfoot
\ExamplesRows
\end{longtable}}

% =================================================================
\chapter{Messages and troubleshooting}\label{ch:messages}

Funis distinguishes three kinds of message. An \emph{input error} means that a field cannot be read. A
message \emph{the task could not be solved} means that the data do not admit a solution of the task, and
gives the reason. A \emph{warning}, shown in colour above the summary, accompanies a solution that holds
but deserves attention. Table~\ref{tab:messages} lists the most frequent ones.

{\small
\begin{longtable}{@{}>{\raggedright\arraybackslash}p{0.44\linewidth}>{\raggedright\arraybackslash}p{0.52\linewidth}@{}}
\caption{Frequent messages and what to do.}\label{tab:messages}\\
\toprule
\bfseries Message & \bfseries Meaning and remedy\\ \midrule
\endfirsthead
\toprule
\bfseries Message & \bfseries Meaning and remedy\\ \midrule
\endhead
\bottomrule
\endfoot
\emph{Row …, column …: '…' is not a number} & A cell of a table holds text. Correct the cell.\\
\emph{The helix needs R > 0 and an end angle beyond the start angle} & Check the radius and the
angles.\\
\emph{The plan of the path stops or reverses} & Two points of the table have the same plan position, or
the plan turns back. Reorder or thin the points.\\
\emph{The path has a stationary elevation} & A supported path needs a non-zero grade, since the thrust
is the vertical force divided by the grade. Give the path a grade everywhere.\\
\emph{The line is not in compression} & The thrust becomes positive somewhere. Raise the preload or
check the loads.\\
\emph{The required reaction leaves the contact cone} & The wall cannot supply the reaction with the
given friction and strip. Read the admissible preloads and multipliers in the summary and in the
\ui{Contact} tab.\\
\emph{The contact cone needs μ > 0 and a positive strip width} & Correct the contact data, or switch
off the contact check.\\
\emph{The assigned path is not compatible with these loads without reactions} & No initial force makes
the path an axial equilibrium. Use a supported path, or find the curve by form-finding.\\
\emph{The horizontal force vanishes} & The curve would turn vertical, where the projected length
stops being a parameter. Change the initial force or the loads.\\
\emph{The endpoint correction stopped before reaching the target} & The second support may be out of
reach. Improve the estimate of $\C$ or check the target.\\
\emph{The horizontal force F₁ is not negative} & The plane line would turn vertical or into tension.
Raise the thrust or change the point of the line.\\
\emph{The point of the line must lie between the supports} & Use $0<x_c<L$.\\
\emph{No thrust line lies within the arch under these loads} & By the safe theorem the arch is not
shown to be safe. Read the geometric safety factor, below 1, and the eccentricities.\\
\emph{The line through the three points lies outside the arch} & The prescribed line is not admissible
at the listed joints. Other lines may be.\\
\emph{The … folds back: the thickness exceeds twice the radius of curvature} & The intrados of a
strongly curved centre line crosses itself. Reduce the thickness or change the shape.\\
\emph{A pointed arch needs a rise larger than half the span} & Raise the rise, or use a circular
arch.\\
\emph{The point load at x = … is not above the extrados of the arch} & Move the load within the span of
the extrados.\\
\emph{Unexpected error} & An error of the program. Save the project and report it with the file on the
issue page of the repository.\\
\end{longtable}}

% =================================================================
\chapter{Hypotheses and limitations}\label{ch:limits}

\begin{itemize}
\item \textbf{Axial equilibrium.} Every task computes equilibrium states of lines carrying an axial force
alone. It does not compute deformations, stresses in a section, buckling or dynamic response. A
horizontal action is a pseudo-static load.
\item \textbf{Spatial paths.} The projected length must be a regular parameter: the plan cannot stop,
turn back or become vertical. A supported path needs a non-zero grade everywhere. The lateral support
supplies horizontal reactions only.
\item \textbf{Contact.} The cone describes the wall as a support with friction and, optionally, a
centred strip. A violated cone is a statement on the support, not on the line, and the admissible
intervals are not collapse loads.
\item \textbf{Form-finding.} The horizontal force must stay positive. A second support is reached by a
local correction from the estimate, which may fail for a target far from it.
\item \textbf{Plane funicular.} The line must remain a function of $x$, with $F_1<0$ everywhere.
\item \textbf{Arches.} The hypotheses of Heyman hold: no tensile strength, unlimited compressive
strength and, with a zero friction coefficient, no sliding. The fill loads the arch but has no strength
and no passive resistance. The arch is cut into a finite number of slices, so the results converge with
the number of slices.
\end{itemize}

% =================================================================
\appendix

\chapter{Quick reference}\label{app:quick}

\section{Keyboard shortcuts}\label{app:keys}

\begin{center}\small
\begin{tabular}{@{}lll@{}}
\toprule
\bfseries Command & \bfseries Windows & \bfseries macOS\\ \midrule
New project & \kbd{Ctrl+N} & \kbd{⌘N}\\
Open & \kbd{Ctrl+O} & \kbd{⌘O}\\
Save & \kbd{Ctrl+S} & \kbd{⌘S}\\
Save as & \kbd{Ctrl+Shift+S} & \kbd{⇧⌘S}\\
Run & \kbd{F5} & \kbd{F5}\\
User guide & \kbd{F1} & \kbd{⌘?}\\
Quit & \kbd{Alt+F4} & \kbd{⌘Q}\\
\bottomrule
\end{tabular}
\end{center}

\section{Command line}\label{app:cli}

\begin{center}\small
\begin{tabularx}{\linewidth}{@{}>{\ttfamily\raggedright\arraybackslash}p{0.42\linewidth}X@{}}
\toprule
\normalfont\bfseries Command & \bfseries Effect\\ \midrule
Funis project.funis & Open the window with a project.\\
Funis -{}-selftest [image.png] & Solve every example without showing the window, save a screenshot and
a report \file{image.log}, and exit with 0 when every example solves (Section~\ref{sec:selftest}).\\
python run\_funis.py [project.funis] & The same from the source code.\\
python -m funis & Start from the source code.\\
\bottomrule
\end{tabularx}
\end{center}

\section{Notation}\label{app:notation}

{\small
\begin{longtable}{@{}>{\raggedright\arraybackslash}p{0.16\linewidth}>{\raggedright\arraybackslash}p{0.8\linewidth}@{}}
\toprule
\bfseries Symbol & \bfseries Meaning\\ \midrule
\endhead
$\e_1,\e_2,\e_3$ & Unit vectors of the axes, $\e_3$ upward\\
$s$, $L$ & Projected arc length and projected length of a spatial path\\
$\x^*(s)$, $\x(s)$, $f(s)$ & Spatial path, its plan and its elevation\\
$\tv$, $\nv$, $\rho$ & Projected unit tangent, left normal and projected curvature\\
$f_{,s}$, $J$ & Grade and ratio of actual to projected length\\
$\uv$, $\T$ & Tangent vector $\tv+f_{,s}\e_3$ and unit tangent of the spatial path\\
$\F$, $S$, $N$, $V$ & Force, thrust, axial force and vertical force\\
$\C$, $\Pv$ & Force at the start and primitive of the loads\\
$\pp$, $\rr$ & Load and reaction per projected metre\\
$w$, $h$, $\beta$, $\gamma$, $\lambda$ & Vertical and horizontal load densities, direction of the
horizontal load, pseudo-static coefficient, multiplier of the horizontal load\\
$T$ & Preload, $T=-S(0)$\\
$r_t$, $r_n$, $r_\nu$ & Tangential, normal and wall-normal reaction\\
$\mu$, $L_{sw}$, $\delta$, $\varphi$, $\eta$ & Friction coefficient, strip width, distance of the strip
from the wall, strip limit and effective coefficient of the cone\\
$\G$, $\R$ & Point force and point reaction\\
$Q$ & Horizontal force of a free curve\\
$\kappa$, $\tau$ & Curvature and torsion of a free curve\\
$B$, $\mv$, $\dv$ & Conserved component, normal and direction of a vertical load plane\\
$x$, $z$, $H$ & Abscissa, height and thrust in the plane tasks\\
$t$, $b$, $e$ & Thickness and depth of an arch, eccentricity of the centre of pressure\\
\bottomrule
\end{longtable}}

% =================================================================
\renewcommand{\bibname}{References}
\begin{thebibliography}{9}
\bibitem{article} C.~Olivieri, \emph{Equilibrium and form-finding of spatial compression paths under
multidirectional loading through Linear Arch Static Analysis}.
\bibitem{funis} C.~Olivieri, \emph{Funis: equilibrium and form-finding of spatial compression paths},
software, Zenodo, \href{https://doi.org/\FunisDOI}{doi:\FunisDOI}.
\bibitem{heyman1966} J.~Heyman, The stone skeleton, \emph{International Journal of Solids and
Structures} 2 (1966) 249--279, \href{https://doi.org/10.1016/0020-7683(66)90018-7}{doi:10.1016/0020-7683(66)90018-7}.
\bibitem{heyman1969} J.~Heyman, The safety of masonry arches, \emph{International Journal of
Mechanical Sciences} 11 (1969) 363--385,
\href{https://doi.org/10.1016/0020-7403(69)90070-8}{doi:10.1016/0020-7403(69)90070-8}.
\end{thebibliography}

\end{document}
